\documentclass[a4paper]{article}

\usepackage[utf8]{inputenc}
\usepackage[T2A]{fontenc}
\usepackage[russian,english]{babel}
\usepackage{amsmath}
\usepackage{booktabs}
\usepackage{siunitx}
\sisetup{
  detect-weight=true,
  detect-inline-weight=math,
  table-number-alignment=center,
}
\setlength{\tabcolsep}{6pt}

\title{Сравнение RRT и GCS: онлайн построение регионов}
\date{}
\begin{document}
\maketitle

Результаты сравнения времени планирования и длины пути для RRT и GCS по парам конфигураций (way: начальная--целевая). RRT выполняется на каждой итерации; GCS --- когда обе конфигурации лежат в построенных IRIS-регионах.

\begin{table}[htbp]
\centering
\caption{Время планирования (с), длина пути и время построения IRIS: RRT vs GCS}
\label{tab:online-gcs-rrt-gcs}
\begin{tabular}{c S[table-format=1.2] S[table-format=1.2] S[table-format=2.2] S[table-format=2.2] S[table-format=2.2]}
\toprule
{Way ($i$--$j$)} & {RRT time (s)} & {GCS time (s)} & {RRT path len.} & {GCS path len.} & {IRIS build (s)} \\
\midrule
$-1$--$0$  & 0.00 & 0.01 & 0.00 & 0.00 & 0.93 \\
$0$--$4$   & 0.01 & 0.31 & 5.48 & 3.38 & 14.06 \\
$0$--$5$   & 1.24 & 0.14 & 15.34 & 8.08 & 0.00 \\
$0$--$6$   & 0.14 & \multicolumn{1}{c}{---} & 7.47 & \multicolumn{1}{c}{---} & 13.46 \\
$1$--$0$   & 0.45 & 0.31 & 18.14 & 8.06 & 0.00 \\
$1$--$6$   & 0.46 & 0.39 & 9.93 & 3.70 & 0.00 \\
$1$--$7$   & 0.14 & 0.53 & 8.97 & 3.23 & 0.00 \\
$2$--$1$   & 0.05 & 0.54 & 8.13 & 2.83 & 0.00 \\
$3$--$3$   & 0.00 & 0.06 & 0.00 & 0.00 & 0.00 \\
$3$--$4$   & 0.01 & 0.39 & 4.98 & 2.45 & 0.00 \\
$3$--$6$   & 0.01 & 0.30 & 6.57 & 1.67 & 0.00 \\
$4$--$1$   & 0.12 & 0.42 & 17.62 & 7.13 & 0.00 \\
$4$--$2$   & 0.02 & 0.53 & 7.11 & 7.26 & 35.89 \\
$4$--$3$   & 0.01 & 0.52 & 4.54 & 2.48 & 0.00 \\
$4$--$7$   & 0.02 & 3.50 & 9.62 & 3.82 & 14.11 \\
$5$--$1$   & 0.04 & 0.19 & 10.55 & 10.14 & 28.26 \\
$5$--$3$   & 0.07 & 0.30 & 9.96 & 5.73 & 15.30 \\
$5$--$6$   & 0.10 & 0.30 & 12.59 & 6.11 & 0.00 \\
$6$--$0$   & 0.03 & 0.12 & 6.62 & 3.38 & 19.42 \\
$6$--$3$   & 0.02 & 0.36 & 4.13 & 1.67 & 0.00 \\
$6$--$4$   & 0.02 & 0.41 & 7.00 & 2.86 & 0.00 \\
$6$--$5$   & 1.72 & 0.25 & 12.55 & 6.11 & 4.34 \\
$7$--$5$   & 1.82 & 0.29 & 16.11 & 4.99 & 0.00 \\
$7$--$6$   & 0.00 & 0.45 & 0.55 & 0.79 & 0.00 \\
\bottomrule
\end{tabular}
\end{table}

Прочерк (---) означает, что для данной пары GCS не использовался. IRIS build (s)~--- время построения IRIS-регионов для данного пути при RRT fallback; 0.00~--- путь планировался по GCS за счёт регионов, построенных на других путях. RRT по времени часто быстрее; GCS даёт более короткие пути.

\medskip

Для сравнения приводим результаты эксперимента \textbf{без умного выбора ключевых точек} (keypoints), с теми же парами конфигураций.

\begin{table}[htbp]
\centering
\caption{Без умного выбора keypoints: время планирования, длина пути и время построения IRIS}
\label{tab:online-gcs-rrt-gcs-nosmart}
\begin{tabular}{c S[table-format=1.2] S[table-format=1.2] S[table-format=2.2] S[table-format=2.2] S[table-format=3.2]}
\toprule
{Way ($i$--$j$)} & {RRT time (s)} & {GCS time (s)} & {RRT path len.} & {GCS path len.} & {IRIS build (s)} \\
\midrule
$-1$--$0$  & 0.00 & 0.01 & 0.00 & 0.00 & 2.04 \\
$0$--$4$   & 0.02 & 1.05 & 5.48 & 3.38 & 51.45 \\
$0$--$5$   & 1.10 & 0.43 & 15.34 & 8.08 & 0.00 \\
$0$--$6$   & 0.15 & \multicolumn{1}{c}{---} & 7.47 & \multicolumn{1}{c}{---} & 132.33 \\
$1$--$0$   & 0.42 & 0.83 & 18.14 & 8.06 & 0.00 \\
$1$--$6$   & 0.43 & 1.41 & 9.93 & 3.70 & 0.00 \\
$1$--$7$   & 0.15 & 1.82 & 8.97 & 3.23 & 0.00 \\
$2$--$1$   & 0.05 & 1.84 & 8.13 & 2.83 & 0.00 \\
$3$--$3$   & 0.00 & 0.25 & 0.00 & 0.00 & 0.00 \\
$3$--$4$   & 0.02 & 1.41 & 4.98 & 2.45 & 0.00 \\
$3$--$6$   & 0.02 & 1.07 & 6.57 & 1.67 & 0.00 \\
$4$--$1$   & 0.06 & 1.47 & 17.62 & 7.13 & 0.00 \\
$4$--$2$   & 0.04 & 1.84 & 7.11 & 7.26 & 90.77 \\
$4$--$3$   & 0.01 & 1.86 & 4.54 & 2.48 & 0.00 \\
$4$--$7$   & 0.02 & 4.56 & 9.62 & 3.82 & 143.09 \\
$5$--$1$   & 0.04 & 0.73 & 10.55 & 10.14 & 126.77 \\
$5$--$3$   & 0.08 & 1.06 & 9.96 & 5.73 & 73.70 \\
$5$--$6$   & 0.11 & 1.36 & 12.59 & 6.11 & 0.00 \\
$6$--$0$   & 0.04 & 0.43 & 6.62 & 3.38 & 100.89 \\
$6$--$3$   & 0.03 & 1.37 & 4.13 & 1.67 & 0.00 \\
$6$--$4$   & 0.04 & 1.63 & 7.00 & 2.86 & 0.00 \\
$6$--$5$   & 1.59 & 1.07 & 12.55 & 6.11 & 48.11 \\
$7$--$5$   & 1.61 & 1.32 & 16.11 & 4.99 & 0.00 \\
$7$--$6$   & 0.00 & 1.78 & 0.55 & 0.79 & 0.00 \\
\bottomrule
\end{tabular}
\end{table}

\medskip

Во втором эксперименте (без smart keypoints) суммарное время IRIS увеличивается (например, для путей $4$--$2$, $4$--$7$, $5$--$1$), что отражает дополнительное построение регионов в «ненужных» местах.

\end{document}
